% Source: https://artofproblemsolving.com/community/c3h2029806_happy_pi_day_maybe_a_competition
% Stable identity: aops-topic-2029806-q15
Number the students \(0,1,\dots,15\) clockwise, starting with the first thrower. Skipping \(x\) students means each toss moves the ball \(x+1\) places, so after \(k\) tosses the ball is at student \(k(x+1)\bmod16\).

The ball reaches every student exactly when the multiples of \(x+1\) cover all residues modulo \(16\), which happens iff \(\gcd(x+1,16)=1\), i.e. \(x+1\) is odd. (If \(d=\gcd(x+1,16)>1\), only multiples of \(d\) are reached.) With \(1\le x+1\le15\), this means \(x\in\{0,2,4,\dots,14\}\), and the sum is
\[
0+2+4+\cdots+14=\boxed{56}.
\]
